%
% komplexitaet.tex -- Beispiel-File für Teil 3
%
% (c) 2026 Nathanael Fässler, OST Ostschweizer Fachhochschule
%
% !TEX root = ../../buch.tex
% !TEX encoding = UTF-8
%

\section{Komplexität}
Dieser Abschnitt setzt theoretische Grundlagen zur \emph{Komplexitätstheorie} voraus, die sehr gut von Müller in seinem Buch~\cite{buchberger:automaten-und-sprachen} erklärt werden.
Bekanntlich ist das Lösen eines Sudokus ein NP-vollständiges Problem.
Als \emph{NP-vollständig} bezeichnet man die Klasse der Probleme die sich von einer nichtdeterministischen Maschine\footnote{Leider existieren deterministische Maschinen nur in der Theorie, 
und daher lassen sich diese Probleme in der Praxis nur in exponentieller Zeit lösen.} 
in polynomieller Zeit lösen lassen.
Es stellt sich die Frage ob das Lösen von Polynomgleichungssystemen auch NP-vollständig ist.

\subsection{Entscheidungsprobleme}
Formell lässt sich Sudoku als das Entscheidungsproblem
\begin{equation}
  \begin{split}
    \text{SUDOKU} = \{ \langle \mathcal{S} \rangle \mid & \mathcal{S} \text{ ist ein Sudoku mit vorgegeben Zahlen, welches eine Lösung besitzt,} \\
    & \text{welches die bekannten Sudoku-Regeln erfüllt.} \}
  \end{split}
  \label{buchberger:sudoku-problem}
\end{equation}
definieren.
Das Entscheidungsproblem beantwortet nur die Frage ob es eine Lösung für das Sudoku gibt oder nicht. Es findet nicht die Lösung selbst.
Es wurde jedoch bewiesen, dass das Entscheiden der Existenz einer Lösung in der gleichen Komplexitätsklasse liegt wie das finden der Lösung.

Für das Problem der Polynomgleichungssysteme wird das Entscheidungsproblem
\begin{equation}
  \begin{split}
    \text{POLYNOM} = \{ \langle \mathcal{P} \rangle \mid & \mathcal{P} \text{ ist ein System von Polynomgleichungen,} \\
    & \text{welches mindestens eine Lösung besitzt.} \}
  \end{split}
  \label{buchberger:polynom-problem}
\end{equation}
definiert.

\subsection{$POLYNOM$ ist NP-vollständig}
\begin{theorem}
  $POLYNOM$ ist NP-vollständig.
\end{theorem}

\subsubsection{Reduktion}
Wie gezeigt, lässt sich $SUDOKU$ durch $POLYNOM$ lösen.
Somit ist $POLYNOM$ mindestens so schwierig wie $SUDOKU$.
Dieser Sachverhalt lässt sich durch die Reduktion
$SUDOKU \leq_{\mathrm{p}} POLYNOM$
ausdrücken.

Dies genügt nicht um zu beweisen, dass $POLYNOM$ NP-vollständig ist.
Es könnte ja sein, dass $POLYNOM$ noch schwieriger ist als NP.

\subsubsection{NP-Zugehörigkeit}

Das Problem $POLYNOM$ lässt sich in polynomieller Zeit verifizieren.
Das heisst es gibt einen Algorithmus, mit dem man in polynomieller Zeit prüfen kann ob eine vorgeschlagene Lösung tatsächlich eine Lösung des Gleichungssystems ist.

Folgender Algorithmus ist ein solcher Entscheider:
\begin{enumerate}
  \item
  Der Algorithmus erhält ein Gleichungssystem und eine Lösung in der Form $x_1 = 5, \dots x_n = 12$.
  \item
  Setze in jeder Gleichung für jede Variable den in der Lösung angegebenen Wert ein.
  \item
  Rechne in jeder Gleichung die Zahlen aus.
  \item
  Überprüfe für jede Gleichung ob der Wert auf der linken und rechten Seite identisch ist.
\end{enumerate}
Die Laufzeit hängt von den Anzahl Gleichungen $n$, den Anzahl Variablen $v$ und der höchsten Potenz $p$ ab.
Schritt 2 besitzt die Laufzeit $O(n \cdot v)$.
Das ausrechnen der Zahlen in Schritt 3 hat Laufzeit $O(n \cdot v \cdot p)$. Das ausrechnen einer Potenz $x^p$ benötigt $p$ Multiplikationen, welche polynomielle Laufzeit haben.
Der letzte Schritt besitzt die Laufzeit $O(n)$.

Somit ergibt sich die Gesamtlaufzeit
\begin{equation}
O(n \cdot v) + O(n \cdot v \cdot p) + O(n) = O(n \cdot v \cdot p)
,
\end{equation}
welche polynomiell ist.

Daraus folgt, dass $POLYNOM$ in NP liegt und nicht ausserhalb.

\begin{proof}
Die NP-vollständigen Probleme sind die Klasse der schwierigsten Probleme innerhalb von NP und sind alle gleich schwierig.
Durch die beiden Resultate
\begin{enumerate}
  \item
  $POLYNOM$ ist mindestens so schwierig wie $SUDOKU$.
  \item
  $POLYNOM$ liegt in NP
\end{enumerate}
lässt sich schlussfolgern, dass $POLYNOM$ gleich schwierig ist wie $SUDOKU$.
Dies lässt sich durch
$SUDOKU \equiv_{\mathrm{P}} POLYNOM$
ausdrücken.
Somit ist bewiesen, dass $POLYNOM$ NP-vollständig ist.
\end{proof}

\subsection{P vs. NP}
Eines der Millennium-Probleme\footnote{Die Milennium-Probleme umfassen eine Liste von ungelösten mathematischen Problemen.
Wer eines davon löst wird mit einer Million US-Dollar belohnt.} 
ist die Frage $P \stackrel{?}{=} NP$, also ob die Komplexitätsklassen P und NP verschieden sind, was von den meisten Mathematikern vermutet wird.
Falls sich herausstellen würde, dass $P = NP$ gilt, würde es bedeuten, dass es für alle NP-vollständigen Probleme einen Algorithmus gibt der das Problem in polynomieller Zeit löst.
Wenn also jemand es schafft einen allgemeinen Algorithmus zu finden der Sudokus oder Polynomgleichungssysteme, welche wie besprochen NP-vollständig sind, löst, hätte das $P = NP$ bewiesen.
Dies hätte schwerwiegende Auswirkungen zum Beispiel auf die Kryptographie:
Die meisten Kryptographischen Systeme basieren auf einer langen Rechenzeit um einen unbekannten \emph{Key} zu berechnen.
Der beste Weg so einen Key zu berechnen ist meistens durch Bruteforcing\footnote{Alle möglichen Varianten ausprobieren.}.
Wenn es auf einmal einen Weg gibt Keys in polynomieller Zeit zu berechnen, wären viele kryptographische Systeme nicht mehr sicher.